\section{Problem and system design}
\label{sec:method}
A collaborative map requires accurate relative placement of robots, beyond accurate independently aligned trajectories. Ground and aerial robots may revisit the same place while observing different surfaces or height ranges. Submaps must therefore contain enough shared spatial structure for retrieval and verification, while remaining economical to exchange and optimize. Our design separates three decisions: where a submap lies, when its geometry is ready to publish, and how its pose is refined in the shared graph.

Each robot runs EllipseLIO~\cite{ellipselio} against its persistent local map. The collaborative backend consumes native retained point representatives, forms spatial submaps, retrieves potential overlaps using MapClosures~\cite{mapclosures}, verifies them with point-cloud registration, and optimizes the resulting submap graph using distributed CBS~\cite{cbs}. Odometry state, covariance, calibration and IMU propagation remain continuous. Circle handovers do not change the odometry map. The backend representation described here is point-cloud based; the ellipsoid representation belongs to the odometry frontend.

\subsection{Incremental placement independent of the trajectory}
Each robot places circles in one fixed gravity-horizontal frame using observed map geometry. The current capture interface supplies chronological 80\,m native area observations, triggered by 20\,m horizontal travel or 10\,s. Their growing union contains unique native point IDs available so far. It is not a complete archive of all raw scans. Robot poses provide coordinate transformations and creation anchors, but are not candidate inputs to the circle-position optimizer.

Let $\mathcal{U}_r(t)$ be that union for robot $r$, and $\Pi_r$ its fixed horizontal projection. Occupied 1\,m grid cells, with fixed origin, are counted once regardless of point density or observation height. Let $\mathcal{G}_r(t)$ denote their centers. A circle has fixed radius $R=40$\,m and center $c_{rk}\in\R^2$. The coverage constraint is
\begin{equation}
 \frac{1}{|\mathcal{G}_r(t)|}\sum_{g\in\mathcal{G}_r(t)}
 \mathbf{1}\!\left[\min_k\|g-c_{rk}\|\le R\right]\ge0.99.
 \label{eq:coverage}
\end{equation}
Existing centers and IDs remain immutable. When new observations lower coverage below the target, deterministic farthest-point selection generates candidate centers from uncovered cells until each is within $R/4$ of a candidate. A sparse candidate-to-cell coverage matrix supports a greedy feasible cover followed by a HiGHS mixed-integer solve minimizing the number of new circles. The solve has a 10\,s time limit and retains the greedy solution if no validated feasible incumbent is available.

Only new centers may be refined continuously. Refinement reduces deviation of each circle's greatest neighbor overlap from a 30\% target while protecting enough covered cells to preserve Eq.~\eqref{eq:coverage}. Disk overlap is intersection area divided by $\pi R^2$; 20--40\% is the desired band, not a hard guarantee. Redundant new circles are removed only after exact coverage validation. Coverage takes priority over circle count, which takes priority over overlap. Any optimality certificate applies to the current candidate set, not to an unknown future map or a globally minimal lifelong cover.

\subsection{Maturing geometry in a fixed footprint}
Footprint coverage does not ensure that the geometry currently inside a circle is useful for place recognition. We therefore create each footprint immediately but postpone its publication. While it is pending, it accumulates every newly available native representative inside that fixed circle, regardless of scan membership. There is no additional vertical crop or scan-age cutoff.

Normal publication requires at least 30\,s of sensor-time age, at most 5\% growth in both unique point IDs and occupied cells over a trailing interval of at least 30\,s, a supported terrain fit, and at least five ORB features summed across the height bands. Circles publish in permanent ID order. A 120\,s sensor-time deadline bounds head-of-line waiting, checked at the next available observation. End-of-stream tails are also released. Deadline and tail releases retain their actual support status; they are not relabeled stable. This minimum-age policy does not promise 30\,s publication latency.

At sealing time $t_s$, the published support is
\begin{equation}
 \mathcal{S}_{rk}=\{p\in\mathcal{U}_r(t_s):\|\Pi_r(p)-c_{rk}\|\le R\}.
\end{equation}
Geometry is expressed in the IMU frame of the original creation scan. The circle center, creation pose and pose timestamp never move to the latest robot position. Separate creation, seal and wall-availability timestamps prevent peers from querying future geometry. Payload hashes bind the descriptor, registration evidence and endpoint metadata. After publication, geometry is immutable; later observations remain available to pending circles and final map visualization but do not silently revise existing factors.

\subsection{Point-cloud BEVs and multilayer matching}
Sealed point clouds are projected onto gravity-horizontal density images at 0.5\,m resolution, using the pinned MapClosures normalization, density threshold 0.05, ORB extractor, Hamming threshold 50 and HBST retrieval. The exchanged representation comprises local ORB descriptors and metric feature coordinates, rather than one global descriptor vector per submap. Candidate similarity supplies a geometric hypothesis, not a guarantee of three-dimensional overlap.

Terrain-relative bands address the height-dependent structures visible to ground and aerial platforms. A spatially balanced low-surface estimator uses the 10th-percentile height of supported 4\,m cells, with at least 20 points per cell. Its plane fit uses 600 seeded hypotheses, 0.35\,m tolerance, at most $12^\circ$ slope, at least 30 supported cells, at least 15\% support, and 8\,m minor-axis spread. Relative height above this plane defines five bands: $[-0.5,2)$, $[2,5)$, $[5,10)$, $[10,20)$ and $[20,\infty)$ metres. This geometric layering is motivated by multi-BEV place recognition~\cite{forestlpr}; it does not use a learned global descriptor. A supported low surface need not be semantic ground. A failed terrain fit disables the layered query; full-height inspection images are not a retrieval fallback in this configuration.

Same-band ORB matches are merged in metric horizontal coordinates, with duplicate endpoint pairs consolidated before joint planar RANSAC. The estimator uses 688 seeded layer-balanced trials, 1.5\,m residual tolerance, one-to-one support within 0.5\,m, and more than five inliers. At most 20 candidate hypotheses are ranked per query; the live verifier admits at most two per query under its queue budget. Same-robot temporal exclusion is 30\,s. These are implementation settings, not accuracy guarantees implied by the referenced methods.

\subsection{Geometric verification and graph factors}
The planar seed is lifted through the gravity transforms, then initialized vertically from the point geometry. The initializer uses 0.8\,m voxels, horizontal neighbors within 1\,m, 0.5\,m height-vote bins and at most five separated modes. It does not use the robot's known flight height or ground truth. Point-cloud GICP~\cite{gicp,smallgicp} refines and verifies the seed using 0.4\,m evidence voxels, a 1.5\,m correspondence radius, 0.6\,m inlier radius, at least 100 inliers, at least 30\% overlap, and at most 0.35\,m inlier RMSE, with observability and conditioning checks.

Verified loops undergo PCM at probability 0.99 with minimum clique size one. A singleton loop may therefore be retained without a second loop corroborating it. Only PCM-retained loops can admit live GPU VGICP factors. The registration adapter uses the native \texttt{gtsam\_points} kernel family, up to 12,000 points per endpoint, 0.5\,m preprocessing, 20 covariance neighbors, and two adaptive voxel levels. The recorded factor maps use 1 and 2\,m voxel resolutions. Initial support checks can reject a geometry factor while preserving its companion verified pose constraint. A conservative information-ratio cap of one limits the registration contribution. Pose and geometry factors reuse observations and are not independent measurements.

\subsection{Persistent distributed optimization and reconstruction}
Each graph variable represents a fixed circle's creation-anchor pose; its ID follows the usual robot/submap endpoint contract. Local odometry supplies relations between creation anchors. Relative-pose loops and admitted geometry factors supply cross-submap constraints. Robot-local CBS optimizers, geometry caches and belief state persist while causal graph prefixes are admitted. Immutable endpoint/session validation and peer exchanges enforce ownership; no centralized optimization is used for the reported solution.

Each robot requests an update every 10\,s. Outgoing pose or anchor exchange is selected randomly each iteration, not in a strict alternating schedule. Both categories update the local graph; they are not separate solver kernels. Anchor beliefs propagate through the connected component, and a common-reference check precedes the final 30 settling iterations. A requested update period is distinct from complete graph-admission, shared-frame publication or convergence latency.

For raw and optimized creation poses $\widehat{X}_{rk}$ and $X_{rk}$, define correction $C_{rk}=X_{rk}\widehat{X}_{rk}^{-1}$. Dense trajectory reconstruction averages corrections sharing one creation timestamp, using mean translation and mean rotation, then interpolates translation and rotation (SLERP) between consecutive timestamps. It applies the resulting correction to every native odometry pose. Map reconstruction instead uses each circle's actual optimized correction: each retained point is assigned to its nearest footprint center and emitted once per robot-local native ID. Points outside all footprints are retained with the nearest correction and flagged. Neither reconstruction introduces a new registration or ground-truth alignment.
